\documentclass{amsart}
\usepackage{amsmath,amssymb,amsthm}

\title{The Braid Permutation $\sigma$ on $\mathbb{Z}/10\mathbb{Z}$}
\author{Brayden Ross Sanders / 7Site LLC}
\date{2026}

\newtheorem{theorem}{Theorem}
\newtheorem{corollary}[theorem]{Corollary}
\newtheorem{remark}[theorem]{Remark}

\begin{document}
\maketitle

\begin{abstract}
The split operator $F$ on $\mathbb{Z}/2\mathbb{Z} \times \mathbb{Z}/5\mathbb{Z}$ induces
a permutation $\sigma$ on the ten operators of $\mathbb{Z}/10\mathbb{Z}$.
We establish that $\sigma$ has exactly four fixed points and one 6-cycle.
\end{abstract}

\section{Statement}

The ten operators of $\mathbb{Z}/10\mathbb{Z}$ carry canonical names:
VOID(0), LATTICE(1), COUNTER(2), PROGRESS(3), COLLAPSE(4),
BALANCE(5), CHAOS(6), HARMONY(7), BREATH(8), RESET(9).

\begin{theorem}[Braid Permutation]
The split operator $F$ on $\mathbb{Z}/2\mathbb{Z} \times \mathbb{Z}/5\mathbb{Z}$
induces the permutation
\[
  \sigma = [0, 7, 1, 3, 2, 4, 5, 6, 8, 9]
  \quad (\text{in one-line notation: } \sigma(i) = \text{table entry at position } i).
\]
This permutation has:
\begin{enumerate}
  \item Exactly 4 fixed points: $\{\mathrm{VOID}(0),\, \mathrm{PROGRESS}(3),\,
        \mathrm{BREATH}(8),\, \mathrm{RESET}(9)\}$.
  \item Exactly one 6-cycle:
  $\mathrm{LATTICE}(1) \to \mathrm{HARMONY}(7) \to \mathrm{CHAOS}(6)
   \to \mathrm{BALANCE}(5) \to \mathrm{COLLAPSE}(4) \to \mathrm{COUNTER}(2) \to \mathrm{LATTICE}(1)$.
\end{enumerate}
The order of $\sigma$ is $\mathrm{lcm}(1,6)=6$, confirmed by $\sigma^6 = \mathrm{id}$.
\end{theorem}

\begin{proof}
Fixed points satisfy $\sigma(i)=i$. By inspection: $\sigma(0)=0$, $\sigma(3)=3$,
$\sigma(8)=8$, $\sigma(9)=9$. All other elements lie in the 6-cycle:
$\sigma(1)=7$, $\sigma(7)=6$, $\sigma(6)=5$, $\sigma(5)=4$, $\sigma(4)=2$,
$\sigma(2)=1$. This is a single cycle of length 6, verified by tracing.
Order $= \mathrm{lcm}(1,6)=6$.
\end{proof}

\begin{remark}[Structural significance]
The fixed points are invariant under the braid dynamics; they are ``inert'' operators.
The 6-cycle elements are ``active'' — they transform under the algebra.
This partition is algebraically forced by the ring structure of $\mathbb{Z}/10\mathbb{Z}$
and carries no free parameters.
\end{remark}

\end{document}
